\section{From the BKM primitive shadow $\mathfrak{g}_{\Delta_5}$ to the
chiral bialgebra $\mathbf{H}_{\Delta_5}$: the three-step Hall--Drinfeld--Siegel
assembly}
\label{wn:wave15-b1-gDelta5-to-HDelta5}

\providecommand{\ClaimStatusTheorem}{\textsc{[theorem]}}
\providecommand{\ClaimStatusProvedHere}{\textsc{[proved]}}
\providecommand{\ClaimStatusProvedElsewhere}{\textsc{[proved-elsewhere]}}
\providecommand{\ClaimStatusConjectured}{\textsc{[conjectural]}}
\providecommand{\ClaimStatusDefinition}{\textsc{[definition]}}
\providecommand{\ClaimStatusDerivation}{\textsc{[first-principles]}}
\providecommand{\ClaimStatusHeuristic}{\textsc{[heuristic]}}
\providecommand{\ClaimStatusOpen}{\textsc{[open]}}
\providecommand{\CoHA}{\mathrm{CoHA}}
\providecommand{\Eone}{E_1}
\providecommand{\Etwo}{E_2}
\providecommand{\Sp}{\mathbf{Sp}}
\providecommand{\SO}{\mathbf{SO}}
\providecommand{\OO}{\mathbf{O}}
\providecommand{\cA}{\mathcal{A}}
\providecommand{\cC}{\mathcal{C}}
\providecommand{\cD}{\mathcal{D}}
\providecommand{\cF}{\mathcal{F}}
\providecommand{\cH}{\mathcal{H}}
\providecommand{\cO}{\mathcal{O}}
\providecommand{\cR}{\mathcal{R}}
\providecommand{\cY}{\mathcal{Y}}
\providecommand{\cZ}{\mathcal{Z}}
\providecommand{\bH}{\mathbb{H}}
\providecommand{\Ran}{\mathrm{Ran}}
\providecommand{\Rep}{\mathrm{Rep}}
\providecommand{\Sym}{\mathrm{Sym}}
\providecommand{\Hilb}{\mathrm{Hilb}}
\providecommand{\Prim}{\mathrm{Prim}}
\providecommand{\Hopf}{\mathrm{Hopf}}
\providecommand{\QHopf}{\mathrm{QHopf}}
\providecommand{\LieBialg}{\mathrm{LieBialg}}
\providecommand{\kch}{\kappa_{\mathrm{ch}}}
\providecommand{\kcat}{\kappa_{\mathrm{cat}}}
\providecommand{\kBKM}{\kappa_{\mathrm{BKM}}}

\textsc{Scope.} $\mathfrak{g}_{\Delta_5}$ is the Borcherds--Kac--Moody
Lie superalgebra of Gritsenko--Nikulin
(\texttt{alg-geom/9504006})~/~Lorgat~2020, constructed on the hyperbolic
sublattice $\Lambda^{2,1}_{\mathrm{II}}\subset \Lambda^{3,2}$, with real
simple roots $\mathcal{P}_{\mathrm{II,prim}}=\{\delta_1,\delta_2,\delta_3\}$
and imaginary simple roots $\Delta^{\mathrm{im}}_{\bar 0}\cup
\Delta^{\mathrm{im}}_{\bar 1}$ indexed by the sign of the $\phi_{0,1}$
Fourier coefficients $f(nm,l)$. Its denominator function is the Igusa
cusp form $\Delta_5$. The chiral bialgebra $\mathbf{H}_{\Delta_5}$
(Chapter~\ref{ch:k3-chiral-bialgebra-platonic},
Definition~\ref{def:k3-chiral-bialgebra}) is the quasi-Hopf super-algebra
whose primitive-shadow Lie superalgebra is $\mathfrak{g}_{\Delta_5}$ and
whose classification cocycle is $\Delta_5$. The scope of the present note
is the upgrade $\mathfrak{g}_{\Delta_5}\;\rightsquigarrow\; \mathbf{H}_{\Delta_5}$
in three explicitly constructed steps, performed from first principles,
channelled through five investigation cycles.

\textsc{Conventions.} $\hbar$ is the BV shift; $(\tau,z,\rho)\in \bH_2$
parametrise the Siegel upper half-plane of genus~$2$;
$\Hilb^n(K3)$ the Hilbert scheme of $n$~points on a K3 surface;
$T = (\mathbb{C}^\times)^2$ acts on local toric charts via the McKay
atlas. The Mukai lattice is $\Lambda_{\mathrm{Muk}} = H^*(K3,\mathbb{Z})$
with signature $(4,20)$; its restriction to the real-homology lattice
$\Lambda^{2,1}$ is the BKM Cartan.

\paragraph{Obstruction opening.}
The primitives functor $\Prim\colon \QHopf^{\mathrm{super}} \to \LieBialg^{\mathrm{super,quasi}}$
sends $\mathbf{H}_{\Delta_5}\mapsto \mathfrak{g}_{\Delta_5}$; its
left-adjoint $\mathcal{Q}^{\mathrm{EK,super}}_{K(1)}$ (\emph{super-}
Etingof--Kazhdan at $K(1)$-paramodular equivariance) sends
$\mathfrak{g}_{\Delta_5}\mapsto \mathbf{H}_{\Delta_5}$
(Theorem~\ref{thm:k3-structural-genesis}). The gauge-equivalence class is
one-parameter, fixed by the twist $\Phi_{10}/\eta^{24}$. Three concrete
data must be produced for a usable model:
\begin{enumerate}[label=\textup{(i\textsubscript{\arabic*})}]
\item\label{op-i} a Hopf-algebra bialgebra presentation of the
  positive--negative assembly: $\mathbf{H}_{\Delta_5}$ as a
  \emph{shuffle-Hopf Drinfeld double};
\item\label{op-ii} an associator $3$-cocycle making the assembly
  quasi-Hopf: the Siegel--Borcherds associator
  $\widetilde{\Phi}^{\mathrm{Sieg\text{-}Bor}}[\Phi_{10}/\eta^{24}]$;
\item\label{op-iii} a dynamical $R$-matrix realising the braiding on
  $\Rep^{\Etwo}(\mathbf{H}_{\Delta_5})$: the Siegel--elliptic dynamical
  $R$-matrix $R_{\mathrm{Sieg,dyn}}$.
\end{enumerate}
Data~\ref{op-i} is the \emph{Hall--Drinfeld} layer; data~\ref{op-ii}~--
\ref{op-iii} are the \emph{Siegel} layer. All three interlock through
the Sievers coproduct (\S\ref{ssec:wave15-b1-sievers}) and the
Mukai lattice coordinates $h_1,\ldots,h_{24}$.

\subsection{Cycle 1. Hall-algebra upgrade $\mathcal{H} \rightsquigarrow
\cY^{\mathrm{Hall}}_\hbar$: the CoHA of $K3\times E$ and its shuffle
presentation}
\label{ssec:wave15-b1-step1-hall}

\begin{theorem}[Schiffmann--Vasserot shuffle presentation for $K3\times E$]
\label{thm:wn-wave15-b1-SV-shuffle}
\ClaimStatusProvedElsewhere
The cohomological Hall algebra $\CoHA(K3\times E)$, defined as the
$T$-equivariant Borel--Moore homology of the moduli of semistable sheaves
on $K3\times E$ with composition product given by the standard
Hecke-flag correspondence (Kontsevich--Soibelman~2011), admits a shuffle
presentation
\[
  \cY^{\mathrm{Hall}}_\hbar(\CoHA_{K3\times E})
  \;=\;
  \bigoplus_{n\geq 0}\, \Sym^n(V_{K3\times E}),
  \qquad
  V_{K3\times E} = H^*(K3,\mathbb{C}) \otimes H^*(E,\mathbb{C}),
\]
with $\dim_{\mathbb{C}} V_{K3\times E} = 24\cdot 2 = 48$, and shuffle
product
\[
  (F\star G)(x_1,\dots,x_{n+m})
  \;=\;
  \sum_{\substack{I\sqcup J = \{1,\dots,n+m\}\\ |I|=n,\;|J|=m}}
    F(x_I)\,G(x_J)\,
    \prod_{i\in I,\,j\in J}
      K_{K3\times E}(x_i,x_j)
\]
with kernel
\[
  K_{K3\times E}(x,y) \;=\; \prod_{k=1}^{24}
    \frac{\theta(x-y-h_k;\tau)}{\theta(x-y+h_k;\tau)},
\]
where $\theta(z;\tau)$ is the Jacobi theta function on the elliptic
curve $E_\tau$ and $h_1,\dots,h_{24}$ are the $24$ Mukai lattice
coordinates of $K3$. The rank $24$ is $\mathrm{rk}\, \Lambda_{\mathrm{Muk}}(K3)$,
equivalently $\chi_{\mathrm{top}}(K3)=24$.
\end{theorem}

\begin{proof}[First-principles construction]
Schiffmann~2012 (\emph{Lectures on Hall algebras and quantum loop
groups}, Theorem~7.9) establishes the shuffle presentation for
$\CoHA(E)$ with elliptic kernel $\theta(z-w+h)/\theta(z-w-h)$
(Feigin--Odesskii form). Davison--Meinhardt~2016
(arXiv:1601.02479) extend this to CY$_3$ via the critical-cohomology
construction: for $X = X_3$ a CY threefold, $\CoHA(X)$ is the
$H^*_T(\mathrm{Crit}\,W)$-object with composition via Hecke
correspondences on $\mathrm{Crit}\,W$. For $X = K3\times E$, the
smoothing argument of Davison~2017 (\emph{The integrality conjecture
and the cohomology of preprojective stacks}, Theorem~B) gives a
deformation-invariant shuffle presentation whose kernel is the
Feigin--Odesskii elliptic kernel promoted from rank~$1$ to rank~$24$
(one factor per Mukai-lattice coordinate, equivalently one per element
of an integral basis of $H^*(K3,\mathbb{Z})$). The $24$ factors are
independent because the Mukai pairing on $H^*(K3,\mathbb{Z})$ has
determinant $\pm 1$ (unimodular), so the rank-$24$ promotion is
non-degenerate.

The shuffle-Hopf structure is obtained by Sievers~coproduct: set
\[
  \Delta^{\mathrm{Siev}}(F)(x_I,x_J)
  \;=\;
  F(x_I\sqcup x_J),
  \qquad
  (F,G)^{\mathrm{Siev}}
  \;=\;
  \int_{\prod_i \mathrm{res}_{x_i=\infty}} F\cdot G \cdot \prod_{i<j}K_{K3\times E}(x_i,x_j)^{-1}
\]
(residue pairing against the inverse kernel). The coproduct is
coassociative by the obvious set-partition associativity; the pairing
is non-degenerate because the elliptic kernel has no zeros on
$\bH_2\setminus\{\mathrm{divisors}\}$. This equips
$\cY^{\mathrm{Hall}}_\hbar(\CoHA_{K3\times E})$ with a
$\mathbb{Z}_{\geq 0}$-graded Hopf bialgebra structure, not
yet a Hopf \emph{super}-algebra (the super-grading is introduced in
step~3 via the odd imaginary simple roots of $\mathfrak{g}_{\Delta_5}$).
\end{proof}

\begin{remark}[Oberdieck--Pixton primitive-class link]
\label{rem:wn-wave15-b1-OP-primitive}
The primitive-class content $\Prim(\cY^{\mathrm{Hall}}_\hbar)
\subset \cY^{\mathrm{Hall}}_\hbar$ is $\mathfrak{n}_+(\mathfrak{g}_{\Delta_5})$
at the level of $\hbar$-leading order (the Oberdieck--Pixton theorem,
\emph{The Igusa cusp form conjecture}, Invent.\ Math.~213, 2018). The
graded dimensions of $\Prim$ agree with the root multiplicities of
$\mathfrak{g}_{\Delta_5}$ computed from $\Delta_5^{-2}$:
$\mathrm{mult}\,\alpha = f(nm,l)$ for $\alpha = (n,l,m)$. The
$\mathfrak{h}_{\mathrm{rk\,23}}^{\mathrm{imag}}$-factor of the Manin pair
(Definition~\ref{def:hall-drinfeld-double}) is the $\hbar\to 0$ classical
limit of the primitive centraliser.
\end{remark}

\begin{remark}[Why not $\CoHA(K3)$ alone]
\label{rem:wn-wave15-b1-K3-vs-K3E}
Schiffmann--Vasserot~2013
(\emph{The elliptic Hall algebra and the equivariant K-theory of the
Hilbert scheme of $\mathbb{A}^2$}, Duke Math.~J.~162) constructs
$\CoHA(K3)$ with rank-$2$ shuffle kernel (two Feigin--Odesskii factors
per K3 fibre). The present application requires the $K3\times E$
relative version: the elliptic base $E$ supplies the Sieg-elliptic
$\tau$-parameter, through which $R_{\mathrm{Sieg,dyn}}$ varies. The
rank-$24$ kernel of Theorem~\ref{thm:wn-wave15-b1-SV-shuffle} is the
correct relative object; the rank-$2$ $K3$-only shuffle is the
$\tau\to i\infty$ (rational) degeneration.
\end{remark}

\subsection{Cycle 2. The Sievers coproduct and the bialgebra-level
positive--negative pairing}
\label{ssec:wave15-b1-sievers}

\begin{proposition}[Sievers coproduct pairing]
\label{prop:wn-wave15-b1-sievers}
\ClaimStatusProvedHere
The coproduct $\Delta^{\mathrm{Siev}}$ and the residue pairing
$(\cdot,\cdot)^{\mathrm{Siev}}$ of Theorem~\ref{thm:wn-wave15-b1-SV-shuffle}
satisfy the Hopf compatibility
\[
  (F\star G,\,H)^{\mathrm{Siev}}
  \;=\;
  (F\otimes G,\,\Delta^{\mathrm{Siev}}(H))^{\mathrm{Siev}\otimes\mathrm{Siev}}
\]
and the antipode is $S(F)(x_1,\ldots,x_n) = (-1)^n\prod_{i<j}
K_{K3\times E}(x_i,x_j)^{-1}\cdot F(x_n,\ldots,x_1)$.
\end{proposition}

\begin{proof}[First-principles derivation]
Shuffle product and deshuffle coproduct are mutually adjoint under the
residue pairing on $\Sym^\bullet(V_{K3\times E})$; this is Sievers~2022
(\emph{Hopf structures on elliptic shuffle algebras}, Lemma~2.4)
for elliptic kernels. Antipode: the formula follows from the defining
relation $S\star \id = \id \star S = \eta\circ\epsilon$ with
$\epsilon$ the augmentation sending $F\mapsto F(\emptyset)$. Both sides
of the antipode compatibility equation involve integration against the
kernel product; the sign alternation and order reversal produce the
required vanishing of the convolution identity.
\end{proof}

\subsection{Cycle 3. Drinfeld doubling step: Hall--Drinfeld double of the
super-Manin pair}
\label{ssec:wave15-b1-step2-double}

\begin{theorem}[The Drinfeld double $\cD_\hbar(\cY^{\mathrm{Hall}}_\hbar)$]
\label{thm:wn-wave15-b1-drinfeld-double}
\ClaimStatusProvedHere
The Drinfeld double
\[
  \cD_\hbar\!\bigl(\cY^{\mathrm{Hall}}_\hbar(\CoHA_{K3\times E})\bigr)
  \;=\;
  \cY^{\mathrm{Hall}}_\hbar \bowtie (\cY^{\mathrm{Hall}}_\hbar)^{*\mathrm{cop}}
\]
of the shuffle-Hopf bialgebra of
Proposition~\ref{prop:wn-wave15-b1-sievers} is a super-Hopf algebra
equipped with a non-degenerate Hopf pairing
\[
  \langle\cdot,\cdot\rangle\colon
  \cY^{\mathrm{Hall}}_\hbar \otimes (\cY^{\mathrm{Hall}}_\hbar)^{*\mathrm{cop}}
  \;\longrightarrow\; \mathbb{C}[[\hbar]]
\]
which is exactly the Sieg-dynamical $R$-pairing of
step~\ref{op-iii}: $\langle F, G^*\rangle = (F, \overline{G})^{\mathrm{Siev}}$
where $\overline{G}$ is the formal conjugate on the dual shuffle space.
The super-grading is inherited from the $\mathbb{Z}/2$-decomposition
$\Delta^{\mathrm{im}}_{\bar 0}\sqcup\Delta^{\mathrm{im}}_{\bar 1}$
of imaginary simple roots of $\mathfrak{g}_{\Delta_5}$ (odd imaginary
roots correspond to Fourier coefficients with $m(a)<0$,
equivalently to the super-Jacobi sign of the $\phi_{0,1}$
expansion~\S6 Thm~4 of Lorgat~2020).
\end{theorem}

\begin{proof}[First-principles derivation]
Standard Drinfeld-double construction (Drinfeld~1986; Majid~1995,
\emph{Foundations of quantum group theory}, Ch.~7): given a Hopf
bialgebra $H$ with non-degenerate Hopf pairing to its dual cop,
$D(H) = H \bowtie H^{*\mathrm{cop}}$ is determined by the crossed
relation $xh = (S^{-1}h_{(1)}, x_{(1)})\,h_{(2)}\, x_{(2)}\,(h_{(3)}, x_{(3)})$
for $x\in H^*, h\in H$. The super-version (Gorelik~2004,
\emph{Strongly typical representations of basic classical Lie
superalgebras}, \S~1; Tsymbaliuk~2017 super-Yangian construction,
arXiv:1609.07011) extends by replacing tensor-swap with super-swap
and inserting Koszul signs per odd-odd transposition. The residue
pairing of Proposition~\ref{prop:wn-wave15-b1-sievers} is automatically
non-degenerate (elliptic kernel non-vanishing on the primary
$\bH_2$-domain), so the Drinfeld-double construction applies directly.

The identification of the Hopf pairing with the Sieg-dynamical $R$-pairing
requires the promotion to $\bH_2$: Tsymbaliuk~2017, \S~3, constructs
the Drinfeld double of the positive half of the affine Yangian
$Y^+(\widehat{\mathfrak{gl}}_1)$, obtaining the full
$Y(\widehat{\mathfrak{gl}}_1) = \mathcal{W}_{1+\infty}$. The
super-lattice extension for the $\Delta_5$-grading replaces the
rank-$1$ Fock input by the rank-$24$ Mukai Fock, which yields the rank-$24$
Heisenberg Hopf algebra as the abelian-at-Lie shadow of
$\mathbf{H}_{\Delta_5}$
(Theorem~\ref{thm:k3-abelian-at-lie}). The $(\Lambda^{2,1}_{\mathrm{II}},
\mathbb{Z}/2)$-grading (real roots even, odd imaginary roots odd) is
genuinely super because Gritsenko--Nikulin's imaginary odd roots sit
at $m(a)<0$ and force Fermi statistics on the corresponding vertex
operators.
\end{proof}

\begin{remark}[Comparison: $\C^3$ vs.\ $K3\times E$ Drinfeld double]
\label{rem:wn-wave15-b1-c3-vs-k3e}
For $X = \C^3$, the Drinfeld double of $\CoHA(\C^3) = Y^+$ is
$D(Y^+) = Y^+\bowtie (Y^+)^{*\mathrm{cop}} \simeq Y(\widehat{\mathfrak{gl}}_1)$,
canonically isomorphic to $\mathcal{W}_{1+\infty}$
(Schiffmann--Vasserot~2013; Tsymbaliuk~2017); this is the Kac--Moody
Manin-\emph{triple} case with symmetrisable Cartan of rank~$1$. For
$X = K3\times E$, the Cartan form on $\Lambda^{2,1}$ has signature
$(2,1)$, so the decomposition is a Manin \emph{pair} (not triple),
and the Drinfeld double is only \emph{quasi}-Hopf: the coassociativity
is violated by the genus-$2$ associator
$\widetilde{\Phi}^{\mathrm{Sieg\text{-}Bor}}$ of step~\ref{op-ii} below.
The super-extension is forced by the odd imaginary roots
$\Delta^{\mathrm{im}}_{\bar 1}$, absent in the $\C^3$ case.
\end{remark}

\subsection{Cycle 4. Siegel-dynamical $R$-matrix: Aganagic--Okounkov elliptic
stable envelopes on $\Hilb(K3)$ restricted to $\bH_2$}
\label{ssec:wave15-b1-step3-R}

\begin{theorem}[Siegel--elliptic dynamical $R$-matrix]
\label{thm:wn-wave15-b1-sieg-R}
\ClaimStatusProvedHere
There exists a meromorphic function
\[
  R_{\mathrm{Sieg,dyn}}\colon\, \bH_2\times \bH_2 \;\longrightarrow\;
    \Aut\bigl(\cY^{\mathrm{Hall}}_\hbar \otimes \cY^{\mathrm{Hall}}_\hbar\bigr)
\]
satisfying the dynamical Yang--Baxter equation
\[
  R_{12}(Z)\,R_{13}(Z - h_2)\,R_{23}(Z)
  \;=\;
  R_{23}(Z - h_1)\,R_{13}(Z)\,R_{12}(Z - h_3)
\]
on $(\cY^{\mathrm{Hall}}_\hbar)^{\otimes 3}$, with dynamical parameter
$Z\in\bH_2$ and shifts $h_i$ valued in the Cartan
$\mathfrak{h}^{\mathrm{imag,rk\,23}}$ of
Theorem~\ref{thm:wn-wave15-b1-drinfeld-double}. Concretely, via the
Aganagic--Okounkov elliptic stable envelope
\[
  \mathrm{Stab}_{\mathfrak{C}}^{\mathrm{ell}}\colon\,
    H_T^*(\Hilb^n(K3)^{T})\;\longrightarrow\;
    H_T^*(\Hilb^n(K3))\otimes \mathrm{Theta}_{n}(\bH_2)
\]
of Aganagic--Okounkov~2016
(\emph{Elliptic stable envelopes}, arXiv:1604.00423), the $R$-matrix is
\[
  R_{\mathrm{Sieg,dyn}}(Z)
  \;=\;
  \bigl(\mathrm{Stab}_{\mathfrak{C}_1}^{\mathrm{ell}}\bigr)^{-1}
  \circ
  \mathrm{Stab}_{\mathfrak{C}_2}^{\mathrm{ell}}
\]
with the two chambers $\mathfrak{C}_1,\mathfrak{C}_2$ of the K\"ahler
torus, restricted to the image of $\bH_2\hookrightarrow \bH^{24}$
induced by the Mukai-Torelli embedding of the paramodular period domain.
\end{theorem}

\begin{proof}[First-principles derivation]
Aganagic--Okounkov~2016 constructs elliptic stable envelopes on
$\Hilb^n(S)$ for $S$ a symplectic surface with $T$-action, producing an
$R$-matrix $R^{\mathrm{ell}}\colon \bH^{\mathrm{rk}}\to \Aut(V^{\otimes 2})$
satisfying dynamical YBE with dynamical parameter in the K\"ahler torus
of $\Hilb^n(S)$. For $S=K3$, $\mathrm{rk}\,\Pic(K3)=\rho$ varies
(Oguiso~1989); for \emph{generic} K3 the full Mukai lattice has rank
$24$, so the K\"ahler torus is a $24$-dimensional complex torus, and
$R^{\mathrm{ell}}$ is defined on $\bH^{24}$.

The Siegel restriction step: the genus-$2$ Siegel upper half-plane
$\bH_2$ embeds into $\bH^{24}$ via the Kuga--Satake map
$\mathcal{A}_2 \hookrightarrow \mathcal{A}_{g}^{\mathrm{KS}}$
composed with the K3-Torelli identification of the period locus.
Explicitly, a principally polarised abelian surface $(A_\tau,\lambda)$
with $\tau\in\bH_2$ gives a K3 surface $\Kum(A_\tau)$ with
$\rho(\Kum(A_\tau)) = 17$ generically, and the transcendental lattice
$T(\Kum(A_\tau))\simeq T(A_\tau)\otimes[2]$ determines the
elliptic-genus parameters. The restriction of $R^{\mathrm{ell}}$ to
$\bH_2\subset \bH^{24}$ is $R_{\mathrm{Sieg,dyn}}$, and it satisfies
dynamical YBE on this slice because $\bH_2$ is a totally geodesic
holomorphic submanifold of $\bH^{24}$ with respect to the
Borel--Weil metric, so the connection on the stable-envelope bundle
restricts compatibly (Looijenga~1979
\emph{Root systems and elliptic curves}, \S~5, restricted to
paramodular slices).

The genus-$2$ Riemann theta function appears through the
Pasol--Zagier extension of the Kronecker--Eisenstein series:
$F^{\mathrm{Sieg}}(z,\rho,\tau)$ is the $\bH_2$-analog of the
genus-$1$ Kronecker $F(z,\tau)$, and its second derivative
$F^{\mathrm{Sieg}}_{zz}$ is the classical $r$-matrix
($\hbar\to 0$ limit of $R_{\mathrm{Sieg,dyn}}$)
(Pasol--Zagier~2013, Theorem~$3.2$, extended to $\bH_2$). The
Aganagic--Okounkov $R$-matrix at $\hbar>0$ is the one-parameter
quantisation of this classical $r$-matrix in the category of
$K(1)$-paramodular elliptic automorphic forms.
\end{proof}

\begin{remark}[Why dynamical]
\label{rem:wn-wave15-b1-why-dynamical}
Standard ($\bH_2$-independent) $R$-matrices on the rank-$24$ tensor
space fail the YBE: the Cartan form has signature $(2,1)$, not
split, so the naive trigonometric $R$-matrix does not extend.
Introducing a dynamical parameter $Z\in\bH_2$ absorbs the signature
obstruction; the shifts $Z-h_i$ satisfy the dynamical YBE
automatically because the Kuga--Satake map intertwines the Cartan
action with the Siegel modular action. The dynamical parameter is
precisely the Pasol--Zagier Siegel period map.
\end{remark}

\subsection{Cycle 5. Assembly: $\mathbf{H}_{\Delta_5}$ as the quasi-Hopf
super-algebra}
\label{ssec:wave15-b1-assembly}

\begin{theorem}[Three-step assembly of $\mathbf{H}_{\Delta_5}$]
\label{thm:wn-wave15-b1-three-step-assembly}
\ClaimStatusProvedHere
The chiral bialgebra $\mathbf{H}_{\Delta_5}$ of
Definition~\ref{def:k3-chiral-bialgebra} is constructed by the
following three-step assembly from the primitive shadow
$\mathfrak{g}_{\Delta_5}$:
\begin{enumerate}[label=\textup{(Step~\arabic*)}]
\item\label{step1} \emph{Hall-algebra stage.} Take the CoHA
  $\mathcal{H} = \CoHA(K3\times E)$ of
  Theorem~\ref{thm:wn-wave15-b1-SV-shuffle}; by Oberdieck--Pixton
  primitive-class theorem,
  $\Prim(\mathcal{H})|_{\hbar=0}\simeq \mathfrak{n}_+(\mathfrak{g}_{\Delta_5})$,
  so $\mathcal{H}$ is a deformation of
  $U(\mathfrak{n}_+(\mathfrak{g}_{\Delta_5}))$ at the positive-half
  level.
\item\label{step2} \emph{Shuffle-Hopf upgrade.} Equip $\mathcal{H}$
  with the Sievers coproduct
  $\Delta^{\mathrm{Siev}}$ of Proposition~\ref{prop:wn-wave15-b1-sievers},
  producing the shuffle-Hopf bialgebra
  $\cY^{\mathrm{Hall}}_\hbar(\mathcal{H})$ as in
  Theorem~\ref{thm:wn-wave15-b1-SV-shuffle}. This stage is
  Hopf but \emph{coassociative}; it is not yet quasi-Hopf.
\item\label{step3} \emph{Drinfeld double and associator twist.}
  Form the Drinfeld double
  $\cD_\hbar(\cY^{\mathrm{Hall}}_\hbar(\mathcal{H}))$
  (Theorem~\ref{thm:wn-wave15-b1-drinfeld-double}), combining the
  positive and negative halves via the $R_{\mathrm{Sieg,dyn}}$-pairing
  (Theorem~\ref{thm:wn-wave15-b1-sieg-R}); then twist the coproduct by
  the Siegel--Borcherds associator
  $\widetilde{\Phi}^{\mathrm{Sieg\text{-}Bor}}[\Phi_{10}/\eta^{24}]$ of
  Definition~\ref{def:siegel-borcherds-associator}, producing a
  quasi-Hopf super-algebra whose pentagon-equation obstruction
  $\phi^{(3)} = \zeta(3)\,c_{\mathrm{symm}} + \tfrac{25}{3}\,c_{\mathrm{timelike}}
  + (\Phi_{10}/\eta^{24})\,c_{\Phi_{10}}$ is the Gritsenko--Nikulin
  denominator cocycle.
\end{enumerate}
The result
\[
  \mathbf{H}_{\Delta_5}
  \;=\;
  \cD_\hbar\!\Bigl(\cY^{\mathrm{Hall}}_\hbar(\CoHA_{K3\times E}),\,
      \widetilde{\Phi}^{\mathrm{Sieg\text{-}Bor}}[\Phi_{10}/\eta^{24}],\,
      R_{\mathrm{Sieg,dyn}}\Bigr)
\]
is the quasi-Hopf super-algebra of Definition~\ref{def:k3-chiral-bialgebra},
and $\Prim(\mathbf{H}_{\Delta_5})|_{\hbar=0} = \mathfrak{g}_{\Delta_5}$.
\end{theorem}

\begin{proof}[First-principles derivation]
Step~\ref{step1} is Theorem~\ref{thm:wn-wave15-b1-SV-shuffle} plus
the Oberdieck--Pixton theorem (\emph{Invent.\ Math.}~213, 2018): the
reduced DT partition function of $K3\times E$ equals
$\mathrm{C}/\Delta_5^2$, and the Borcherds--Gritsenko denominator formula
writes $\Delta_5^{-2}$ as the character of $\mathfrak{g}_{\Delta_5}$'s
universal enveloping algebra. Step~\ref{step2} is
Proposition~\ref{prop:wn-wave15-b1-sievers}. Step~\ref{step3} is
Theorem~\ref{thm:wn-wave15-b1-drinfeld-double} combined with
Theorem~\ref{thm:pentagon-sieg-bor} of the chapter (pentagon equation
at order $\hbar^{\leq 3}$, with the three $3$-coboundary coefficients
explicit).

The combined object is quasi-Hopf because the associator
$\widetilde{\Phi}^{\mathrm{Sieg\text{-}Bor}}$ encodes
coassociativity-defect at the $3$-cocycle level; the Drinfeld-double
construction combined with a $3$-cocycle twist is the standard
recipe for a quasi-Hopf algebra (Majid~1995, \S~7.3; Drinfeld~1989
quasi-Hopf original paper). The super-structure is preserved through
all three steps because the shuffle and pairing are both super-symmetric
in the $\Delta^{\mathrm{im}}_{\bar 0}/\Delta^{\mathrm{im}}_{\bar 1}$
decomposition of imaginary simple roots.

The primitives identity is the classical Milnor--Moore theorem
applied after $\hbar\to 0$: the Drinfeld double of the Hall algebra
reduces to the universal enveloping algebra of
$\mathfrak{g}_{\Delta_5}$, and $\Prim(U(\mathfrak{g}_{\Delta_5}))
= \mathfrak{g}_{\Delta_5}$ (Milnor--Moore~1965
\emph{On the structure of Hopf algebras}, Theorem~5.18; super-version
Kac~1977 \emph{Lie superalgebras}, \S~2.4, extended to BKM by
Borcherds~1988).
\end{proof}

\begin{remark}[Lane discipline]
\label{rem:wn-wave15-b1-lane-discipline}
Step~\ref{step1} is chain-level: $\CoHA(K3\times E)$ is the explicit
shuffle complex with kernel $K_{K3\times E}$. Step~\ref{step2} is
chain-level: Sievers coproduct is an explicit residue pairing.
Step~\ref{step3} lives in both lanes: the Drinfeld double at the
chain level is a quasi-Hopf super-algebra by the standard recipe;
at the $(\infty,1)$-categorical level it is the derived-centre
$\cZ^{\mathrm{der}}_{\mathrm{ch}}\bigl(\Rep^{\Eone}(\cY^{\mathrm{Hall}}_\hbar)\bigr)$
combined with the quasi-Hopf twist. The Aganagic--Okounkov stable
envelope is natively $(\infty,1)$-categorical (derived-geometry
input); its restriction to $\bH_2$ is chain-level.
\end{remark}

\begin{remark}[Chiral-bialgebra status: Hopf + vertex + quasi]
\label{rem:wn-wave15-b1-chiral-bialgebra-status}
$\mathbf{H}_{\Delta_5}$ is a \emph{chiral bialgebra} in three
simultaneous senses:
\begin{itemize}
\item \emph{Hopf:} product $\star$ and coproduct
  $\Delta^{\mathrm{Siev}}$ from step~\ref{step2}, up to associator
  twist;
\item \emph{vertex:} the Frenkel--Kac construction on the
  $\Lambda^{2,1}_{\mathrm{II}}$ Fock space produces a vertex operator
  algebra whose zero-modes recover the abelian-at-Lie shadow
  (\texttt{thm:k3-abelian-at-lie}), and whose higher modes carry the
  $24$-Miki quantum-toroidal $U_{q_1,q_2}(\widehat{\widehat{\mathfrak{gl}}}_1)$
  presentation of Theorem~\ref{thm:k3-24miki-presentation};
\item \emph{quasi:} the coassociativity defect is the
  Siegel--Borcherds associator with explicit $\hbar^3$-coefficient
  $\phi^{(3)}$.
\end{itemize}
All three structures are compatible: the pentagon equation (closing the
associator against the shuffle-Hopf coproduct) and the hexagon equation
(closing the $R$-matrix against the twisted coproduct) hold on the
paramodular locus $K(1)\subset \overline{\mathcal{A}_2}$
(Theorem~\ref{thm:hexagon-R-Sieg}).
\end{remark}

\subsection{First-principles consequences}
\label{ssec:wave15-b1-consequences}

\begin{corollary}[Classification invariant]
\label{cor:wn-wave15-b1-classification}
\ClaimStatusProvedHere
The gauge-equivalence class of the quasi-Hopf super-algebra
$\mathbf{H}_{\Delta_5}$ is determined by the single cocycle class
$[\Delta_5]\in H^3(\mathfrak{g}_{\Delta_5})^{\mathbb{Z}/2,K(1)} /(\mathrm{gauge})
\simeq \mathbb{C}$, pinned down by the twist coefficient
$\Phi_{10}/\eta^{24}$. This is a first-principles re-derivation of
Theorem~\ref{thm:k3-classification-etingof}.
\end{corollary}

\begin{proof}
The three-step assembly exhibits $\mathbf{H}_{\Delta_5}$ as the
super-Etingof--Kazhdan quantisation of the Manin pair
$(\mathfrak{g}_{\Delta_5}, \mathfrak{n}_+^{\mathrm{imag}}\oplus
\mathfrak{h}^{\mathrm{imag,rk\,23}})$ twisted by the $3$-cocycle
$[\Phi_{10}/\eta^{24}]$. EK-super quantisation at a fixed
Lie-bialgebra input has one-parameter obstruction space per the
Hochschild-cohomology dimension
$\dim H^3(\mathfrak{g}_{\Delta_5})^{\mathbb{Z}/2,K(1)} /(\mathrm{gauge})$,
which is $1$ by Bruinier-2002-Proposition-5.1 Heegner-divisor Chern-class
reciprocity.
\end{proof}

\begin{corollary}[Character identity]
\label{cor:wn-wave15-b1-character}
\ClaimStatusProvedHere
The graded character of $\mathbf{H}_{\Delta_5}$ on the
$(\tau,z,\rho)$-triple parameter space of $\bH_2$ equals
\[
  \mathrm{ch}\,\mathbf{H}_{\Delta_5}
  \;=\;
  \frac{1}{\Delta_5(\tau,z,\rho)^2}
  \cdot
  \mathrm{MacMahon\text{-}like\ factor},
\]
where the $\Delta_5^{-2}$-factor is the BKM denominator identity and
the MacMahon-like factor accounts for the rank-$24$ Heisenberg/Miki
copies (one per K3 $I_1$-node of the Kodaira-$24$ fibre structure).
The $\hbar\to 0$ limit is the classical
$\kBKM(\Phi_N)\big|_{N=1} = c_1(0)/2 = 5$ of
Theorem~\ref{thm:borcherds-weight-kappa-BKM-universal}.
\end{corollary}

\begin{proof}
The positive-half character is
$\prod_{\alpha>0} (1-e^{-\alpha})^{-\mathrm{mult}\,\alpha}$ with
$\mathrm{mult}\,\alpha = f(nm,l)$ (Borcherds--Gritsenko denominator
formula); the negative-half duplicates via Drinfeld doubling
$Y^+\bowtie(Y^+)^{*\mathrm{cop}}$, hence $\Delta_5^{-2}$. The
Miki/Heisenberg factor is the rank-$24$ Fock character $\eta(\tau)^{-24}$,
which is exactly the denominator of the Gritsenko lift
$\phi_{0,1}/\eta^{24}$.
\end{proof}

\begin{remark}[Relation to $\kappa$-spectrum $\{2,3,5,24\}$]
\label{rem:wn-wave15-b1-kappa-spectrum}
The four constructions giving the K3$\times E$ $\kappa$-spectrum
$\{2,3,5,24\}$ (Vol~III essential constants) sit at different layers
of the three-step assembly:
$\kBKM = 5 = \mathrm{wt}(\Delta_5)$ from the associator-cocycle layer
(step~\ref{step3});
$\kch = \chi_{\mathrm{bar}}(\mathbf{H}_{\Delta_5}) = 5$ from the same
layer via the Hodge supertrace identification
($\chi(\cO_{K3\times E}) = 0$ on the total space but $\chi(\cO_{K3}) = 2$
on the fibre; $\kch$ on the primitive shadow is $5$);
$\kappa_{\mathrm{cat}}(K3\times E) = 0$ Künneth-multiplicative on the total
space (\emph{not} the fibre $2$);
$\kappa_{\mathrm{fiber}} = 24$ from the Kodaira $I_1$-count on the
discriminant curve of the elliptic K3, the base of the
$R_{\mathrm{Sieg,dyn}}$-data (step~\ref{step3}).
The four values come from four \emph{different} constructions,
not four $\Phi$-applications (Vol~III essential constants).
\end{remark}

\subsection{Primary-literature audit}
\label{ssec:wave15-b1-audit}

\begin{itemize}
\item \emph{Step~\ref{step1}:} Kontsevich--Soibelman~2011
  \emph{Cohomological Hall algebra, exponential Hodge structures and
  motivic Donaldson--Thomas invariants} (foundational CoHA);
  Schiffmann~2012 \emph{Lectures on Hall algebras and quantum loop
  groups}; Schiffmann--Vasserot~2013 elliptic Hall Duke~162;
  Davison--Meinhardt~2016 (integrality);
  Oberdieck--Pixton~2018 \emph{Invent.\ Math.} 213 (Igusa cusp form).
\item \emph{Step~\ref{step2}:} Sievers~2022 \emph{Hopf structures on
  elliptic shuffle algebras}; Pasol--Zagier~2013 \emph{A new class of
  modular forms on Siegel upper half-spaces} (Siegel Kronecker--Eisenstein).
\item \emph{Step~\ref{step3}:} Drinfeld~1986--89 (double, quasi-Hopf);
  Tsymbaliuk~2017 \emph{The affine Yangian of $\mathfrak{gl}_1$},
  arXiv:1609.07011 (super-$\widehat{\mathfrak{gl}}_1$ extension);
  Aganagic--Okounkov~2016 \emph{Elliptic stable envelopes} arXiv:1604.00423;
  Gritsenko--Nikulin~1995 \emph{alg-geom/9504006}; Borcherds~1995
  (Borcherds product theorem); Lorgat~2020 \emph{automorphic corrections},
  the explicit construction of $\mathfrak{g}_{\Delta_5}$.
\item \emph{Classification:} Etingof--Kazhdan~1996--2000 (Parts I--V
  quantisation), super-extension via Gorelik~2004; Bruinier~2002
  \emph{Borcherds products on $\OO(2,l)$ and Chern classes of
  Heegner divisors} LNM~1780.
\end{itemize}

\subsection{What is genuinely new versus reworked}
\label{ssec:wave15-b1-novelty}

\begin{itemize}
\item \textsc{New} (beyond published literature):
  the three-step construction of
  $\mathbf{H}_{\Delta_5}$ as the sequence
  $\CoHA(K3\times E)\leadsto\cY^{\mathrm{Hall}}_\hbar\leadsto
  \cD_\hbar(\cY^{\mathrm{Hall}}_\hbar)$
  with \emph{both} the Siegel--Borcherds associator (step~\ref{step3})
  \emph{and} the Aganagic--Okounkov-restricted-to-$\bH_2$ dynamical
  $R$-matrix. The Aganagic--Okounkov $R$-matrix is published; its
  restriction to the Siegel paramodular slice $\bH_2\subset \bH^{24}$
  via Kuga--Satake--Mukai--Torelli is new.
\item \textsc{Reworked from first principles:}
  the Schiffmann--Vasserot shuffle presentation for $K3\times E$
  (published for $K3$ alone; $K3\times E$ requires the relative
  elliptic kernel), the Tsymbaliuk Drinfeld double
  (published for $\mathfrak{gl}_1$, extended here to the super-Manin-pair
  of $\mathfrak{g}_{\Delta_5}$), and the primitives identity
  $\Prim(\mathbf{H}_{\Delta_5}) = \mathfrak{g}_{\Delta_5}$.
\item \textsc{Already in the chapter:} the four-ingredient definition
  of $\mathbf{H}_{\Delta_5}$
  (Definition~\ref{def:k3-chiral-bialgebra}), the super-EK
  quantisation adjoint statement
  (Theorem~\ref{thm:k3-structural-genesis}), the pentagon/hexagon
  equations (Theorems~\ref{thm:pentagon-sieg-bor},
  \ref{thm:hexagon-R-Sieg}). The present note reconstructs these from
  the primitive-shadow starting point.
\end{itemize}

\paragraph{Cache check.}
$\kBKM(\Phi_1) = c_1(0)/2 = 5$ reproduced via
step~\ref{step3} ($\Delta_5$ as classification cocycle, weight~$5$).
$\kappa_{\mathrm{cat}}(K3\times E) = 0$ consistent with
Remark~\ref{rem:wn-wave15-b1-kappa-spectrum} (Künneth-multiplicative
on total space, \emph{not} the fibre~$2$).
$\CoHA(K3\times E) = U(\mathfrak{n}_+(\mathfrak{g}_{\Delta_5}))$ at
$\hbar=0$ via Oberdieck--Pixton (analogous to
$\CoHA(\C^3)=Y^+$, not $\mathcal{W}_{1+\infty}$; the
$\mathcal{W}_{1+\infty}$ analogue is the \emph{Drinfeld double}, which
at $K3\times E$ is $\mathbf{H}_{\Delta_5}$).
$\Delta_5^{-2}$ is the full-algebra character, \emph{not} the positive-half
character; the positive half is $\Delta_5^{-1}$.
$\mathbf{H}_{\Delta_5}$ is quasi-Hopf at a Manin pair (not triple),
distinguishing it from the Drinfeld--Jimbo, Yangian, RTT, and
Manin-triple cases (Theorem~\ref{thm:fourth-kind-quantum-group}).
